\chapter{Lower bounds for General Formulas}\label{chap:gen-formulas}

\section{Kalorkoti's Lower Bound}\label{sec:Kalorkoti}

This section would be devoted to the proof of Kalorkoti's lower bound \cite{k85} for formulas computing $\Det_n$ or $\Perm_n$.

\begin{theorem}[\cite{k85}]\label{thm:kalorkoti}
  Any arithmetic formula computing $\Perm_n$ (or $\Det_n$) requires $\Omega(n^3)$ size.
\end{theorem}

Although this is an $\Omega(n^{3/2})$ lower bound, where $n$ is the number of variables, the same technique can also give an $\Omega(n^2)$ lower bound for the following polynomial.

\begin{theorem}[\cite{k85}]\label{thm:kalorkoti-b}
  Any arithmetic formula computing $\sum_{i=1}^n \sum_{j=1}^n x_i^j y_j$ requires $\Omega(n^2)$ size.
\end{theorem}

In this section, we shall prove the lower bound for $\Det_n$ and $\Perm_n$ and leave \autoref{thm:kalorkoti-b} as an exercise. 

The exploitable weakness in this setting is again to use the fact that the polynomials computed at intermediate gates share many polynomial dependencies.

\begin{definition}[Algebraic independence]
  A set of polynomials $\inbrace{f_1,\dots, f_m}$ is said to be \emph{algebraically independent} if there is no non-trivial polynomial $H(z_1,\dots, z_m)$ such that $H(f_1,\dots, f_m)=0$.

  The size of the largest algebraically independent subset of $\vecf=\inbrace{f_1,\dots, f_m}$ is called the \emph{transcendence degree} (denoted by $\mathrm{trdeg}(f)$).
\end{definition}
The proof of Kalorkoti's theorem proceeds by defining a \emph{complexity measure} using the above notion of algebraic independence.
\\


{\bf The Measure:} 
For any subset of variables $Y\subseteq X$, we can write a polynomial $f \in \F[X]$ of the form $f = \sum_{i=1}^s f_i \cdot m_i$ where $m_i$'s are distinct monomials in the variables in $Y$, and $f_i \in F[X \setminus Y]$.
We shall denote by $\mathrm{td}_Y(f)$ the transcendence degree of $\inbrace{f_1,\dots, f_s}$


Fix a partition of variables $X = X_1 \sqcup \dots \sqcup X_r$.
For any polynomial $f\in \F[X]$, define the map $\CM{Kal}:\F[X]\rightarrow \Z_{\geq 0}$ as
$$
\CM{Kal}(f) \quad=\quad \sum_{i=1}^r \mathrm{td}_{X_i}(f).
$$

The lower bound proceeds in two natural steps:
\begin{enumerate}
\item Show that $\CM{Kal}(f)$ is \emph{small} whenever $f$ is computable by a \emph{small} formula. 
\item Show that $\CM{Kal}(\Det_n)$ is \emph{large}. 
\end{enumerate}

\subsection{Upper bounding $\CM{Kal}$ for a formula}

\begin{lemma}\label{lem:kal-upperbound}
  Let $f$ be computed by a fan-in two formula $\Phi$ of size $s$.
Then for any partition of variables $X = X_1\sqcup \dots \sqcup X_r$, we have $\CM{Kal}(f) = O(s)$.
\end{lemma}
\begin{proof}
  For any node $v\in \Phi$, let $\textsc{Leaf}(v)$ denote the leaves of the subtree rooted at $v$ and let $\textsc{Leaf}_{X_i}(v)$ denote the leaves of the subtree rooted at $v$ that are in the set $X_i$.
Since the underlying graph of $\Phi$ is a tree, it follows that the size of $\Phi$ is bounded by a twice the number of leaves.
For each set $X_i$, we shall show that $\mathrm{td}_{X_i}(f) = O(\abs{\textsc{Leaf}_{X_i}(\Phi)})$, which would prove the required bound.
\\

Fix an arbitrary set $Y = X_i$.
Since the goal is to just get a bound on $\textsc{Leaf}_{Y}(\Phi)$, we shall modify the formula $\Phi$ by introducing new leaf variables that compute polynomial in $\F[X \setminus Y]$.
This does not affect the size of $\textsc{Leaf}_{Y}(\Phi)$.
This in some sense is allowing $\Phi$ to ``freely'' compute any polynomial on variables outside $Y$, as we are only interested in how many times the $Y$ variables are used in the computation of $\Phi$.

We modify $\Phi$ by introducing a new type of gate $\boxtimes$ that takes a node $g$ and two leaves $\ell_1$ and $\ell_2$ and computes $\boxtimes(g,\ell_1,\ell_0) = (\ell_1 \cdot g) + \ell_0$.
We shall always ensure that when we introduce new leaf nodes, they only hold polynomials $\F[X \setminus Y]$.

Define the following three sets of nodes:
  \begin{eqnarray*}
    V_0 & = & \setdef{v\in \Phi}{\abs{\textsc{Leaf}_{Y}(v)} = 0}\\
    V_1 & = & \setdef{v\in \Phi}{\abs{\textsc{Leaf}_{Y}(v)} = 1 \quad\text{and}\quad \abs{\textsc{Leaf}_{Y}(\textsc{Parent}(v))} \geq 2}\\
    V_2 & = & \setdef{v\in \Phi}{\abs{\textsc{Leaf}_{Y}(v)} \geq 2}.
  \end{eqnarray*}


  Each node $v\in V_0$ computes a polynomial in $f_v \in \F[X\setminus Y]$, and we shall replace the subtree at $v$ by a leaf computing the polynomial $f_v$.

  Similarly, any node $v\in V_1$ computes a polynomial of the form $f_1 \cdot y_v + f_0$ for some $y_v\in Y$
  and $f_0,f_1 \in \F[X\setminus Y]$. 
  We shall create two new leaf nodes $\ell_0,\ell_1$ computing $f_0$ and $f_1$ respectively, and replace the gate $v$ by a $\boxtimes$ gate with inputs $y_v,\ell_1,\ell_0$ so that it computes $\boxtimes(y_v,\ell_1,\ell_0) = (f_1 \cdot y_v) + f_0$

  Hence, the formula $\Phi$ now reduces to a smaller formula $\Phi_Y$ with
  leaves  being the nodes in $V_0$ and $V_1$ (and nodes in $V_2$ are
  unaffected). 
  Furthermore, all new leaves that are added compute polynomials in $\F[X \setminus Y]$ and hence $\textsc{Leaf}_Y$ is unchanged. 
We would like to show that the size of the reduced
  formula, which is at most twice the number of its leaves, is
  $O(\abs{\textsc{Leaf}_Y(\Phi)})$.

  \begin{observation}\label{obs:v1-bound}$\abs{V_1} \leq \abs{\textsc{Leaf}_Y(\Phi)}$.    
  \end{observation}
  \begin{myproof}{Obs}
    Each node in $V_1$ has a distinct leaf labelled with a variable in $Y$.
Hence, $\abs{V_1}$ is bounded by the number of leaves labelled with a variable in $Y$.
  \end{myproof}
  
  This shows that the $V_1$ leaves are not too many.
Unfortunately, we cannot immediately bound the number of $V_0$ leaves, since we could have a long chain of $V_2$ nodes each with one sibling being a $V_0$ leaf.
The following observation would show how we can eliminate such long chains.

  \begin{observation}\label{obs:same-leaf-collapse}
    Let $u$ be an arbitrary node, and $v$ be another node in the subtree rooted at $u$ with $\textsc{Leaf}_Y(u) = \textsc{Leaf}_Y(v)$.
Then the polynomial $g_u$ computed at $u$ and the polynomial $g_v$ computed at $v$ are related, i.e., $g_u = f_1 g_v + f_0$ for some $f_1,f_0 \in \F[X\setminus Y]$.
  \end{observation}
  \begin{myproof}{Obs}
    If $\textsc{Leaf}_Y(u) =\textsc{Leaf}_Y(v)$, then every node on
    the path from $u$ to $v$ must have a $V_0$ leaf as the other child. 
The
    observation follows as all these nodes are $+$ or $\times$ gates.
  \end{myproof}

  Using the above observation, we shall remove the need for $V_0$
  nodes completely by adding new $\F[X\setminus Y]$ leaves and $\boxtimes$ gates. 
  Formally, if a $V_0$ node computing $f$ was connected to a $+$ node that was computing $f + g$, then we can replace the $V_0$ node by a new leaf $\ell_f$ computing $f$ and replace that $+$ node with $\boxtimes(g,1,\ell_f) = (g \cdot 1) + f$. Similarly, if a $V_0$ node computing $f$ was incident on a $\times$ node that was computing $f \cdot g$, then we can replace the $V_0$ node by a new leaf $\ell_f$ computing $f$ and replace the $\times$ gate with a $\boxtimes(g,\ell_f,0) = g \cdot f + 0$. 

Furthermore, using \autoref{obs:same-leaf-collapse}, we can now contract any two nodes $u$ and $v$ with $\textsc{Leaf}_Y(u) = \textsc{Leaf}_Y(v)$, we can replace the entire chain from $u$ to $v$ by an appropriate $\boxtimes$ node. 
This ensures that no $\boxtimes$ node is connected to another $\boxtimes$ node. 
Overall, $\Phi$ has been transformed in the following ways:
\begin{enumerate}
\itemsep 0pt
\item all $V_0$ nodes are replaced by leaves,
\item all $V_1$ are replaced by $\boxtimes$ nodes with three leaves incident on it,
\item no $\boxtimes$ node is incident on another $\boxtimes$ node, and hence its parent is a $V_2$ node,
\item distinct vertices $u,v\in V_2$ satisfy $\textsc{Leaf}_Y(u) \neq \textsc{Leaf}_Y(v)$, 
\item leaves in $Y$ are untouched,
\item all new leaves compute polynomials only in $\F[X \setminus Y]$.
\end{enumerate}
(2) implies that the number of $|V_1|$ nodes in $\hat{\Phi}_Y$ is at most $\abs{\textsc{Leaf}_Y(\Phi)}$.
Also (4) implies the number of $V_2$ nodes is at most $\abs{\textsc{Leaf}_Y(\Phi)} - 1$.
Therefore, the size of the augmented formula $\hat{\Phi}_Y$ is at most $2\abs{\textsc{Leaf}_Y(\Phi)}$.
\\

Suppose $\Phi$ computes a polynomial $f$, which can be written as $f = \sum_{i=1}^t f_i\cdot m_i$ with $f_i \in \F[X\setminus Y]$ and $m_i$'s being distinct monomials in $Y$.
Since $\hat{\Phi}_Y$ also computes $f$, each $f_i$ is a polynomial combination of the leaves of $\hat{\Phi}_Y$ that compute polynomials in $\F[X \setminus Y]$.
Since $\hat{\Phi}_Y$ consists of at most $2\abs{\textsc{Leaf}_Y(\Phi)}$ augmented nodes, we have that $\mathrm{td}_Y(f) \leq 4\abs{\textsc{Leaf}_Y(\Phi)}$.
Therefore,
$$
\mathrm{td}_Y(f) \quad=\quad \mathrm{trdeg}\setdef{f_i}{i\in [t]} \quad\leq\quad 4\abs{\textsc{Leaf}_Y(\Phi)}
$$
Hence, 
$$
\CM{Kal}(\Phi) = \sum_{i=1}^r \mathrm{td}_{X_i}(f_i) \leq 4\inparen{\sum_{i=1}^r \abs{\textsc{Leaf}_{X_i}}} = O(s).\qedhere
$$
\end{proof}

\subsection{Lower bounding $\CM{Kal}(\Det_n)$}


\begin{lemma}\label{lem:kal-lowerbound}
  Let $X = X_1 \sqcup \dots \sqcup X_n$ be the partition as defined by
  $X_t = \setdef{x_{ij}}{i-j\equiv t\bmod{n}}$. 
Then,
  $\CM{Kal}(\Det_n) = \Omega(n^3)$.
\end{lemma}
\begin{proof}
  By symmetry, it is easy to see that $\text{td}_{X_i}(\Det_n)$ is
  the same for all $i$. 
Hence, it suffices to show that
  $\text{td}_{Y}(\Det_n) = \Omega(n^2)$ for $Y = X_n = \inbrace{x_{11},\dots, x_{nn}}$. 

  To see this, observe that the determinant consists of the monomials
  $\inparen{\frac{x_{11}\dots x_{nn}}{x_{ii}x_{jj}}}\cdot
  x_{ij}x_{ji}$ for every $i\neq j$. 
Hence, $\text{td}_{Y}(\Det_n)
  \geq \mathrm{trdeg}\setdef{x_{ij}x_{ji}}{i\neq j} =
  \Omega(n^2)$. 
Therefore, $\CM{Kal}(\Det_n) =
  \Omega(n^3)$.
\end{proof}

The proof of \autoref{thm:kalorkoti} follows from
\autoref{lem:kal-upperbound} and \autoref{lem:kal-lowerbound}.\\

\begin{exercise}
Prove an $\Omega(n^2)$ lower bound for $\CM{Kal}(f)$ where $f = \sum_{i=1}^n \sum_{j=1}^n x_i^j y_j$ for an appropriate partition. 
\end{exercise}



%%% Local Variables: 
%%% mode: latex
%%% TeX-master: "main"
%%% End: 
